\documentclass[../../main.tex]{subfiles}

\begin{document}

\section{非负可测函数的 Lebesgue 积分}\label{section:ch04-非负可测积分}

$E$ 上的可测函数 $f$ 称为在有限测度集外消失, 若存在 $E$ 的子集 $E_{0}$, 使得 $m(E_{0})<\infty$ 且在 $E\sim E_{0}$ 上 $f\equiv 0$. 说在有限测度集外消失的可测函数 $f$ 有有限支撑, 且定义它的支撑为 $\{x\in E\mid f(x)\neq 0\}$ 是方便的. 在前一节, 我们定义有界可测函数 $f$ 在有限测度集 $E$ 上的积分. 然而, 即便 $m(E)=\infty$, 若 $f$ 是在 $E$ 上的有界可测函数但具有有限支撑, 我们能定义它在 $E$ 上的积分为
\begin{align*}
\int_{E}f=\int_{E_{0}}f,
\end{align*}
其中 $E_{0}$ 有有限测度而在 $E\sim E_{0}$ 上 $f\equiv 0$. 该积分是正确定义的, 即它与使得 $f$ 在其外消失的有限测度集 $E_{0}$ 的选取无关.

\begin{definition}[非负可测函数的积分]\label{definition:ch04-非负积分}
对于 $E$ 上的非负可测函数 $f$, 我们定义 $f$ 在 $E$ 上的积分为
\begin{align}
\int_{E}f=\sup\left\{\int_{E}h\ \middle|\ h\text{ 有界、可测、支撑有限且在 }E\text{ 上 }0\leqslant h\leqslant f\right\}.\label{align:ch04-非负积分定义}
\end{align}
\end{definition}

这是非负扩充实值可测函数的积分的定义, 但它不是意味着这样的函数是可积的定义. 积分是非负的, 这是由于它定义为非负数集的上确界. 但积分可能等于 $\infty$, 比方说, 对在 $E$ 的无穷测度子集上取正常数值或在 $E$ 的正测度的子集上取 $\infty$ 值的非负可测函数, 就出现这种情况.

\begin{proposition}[Chebyshev 不等式]\label{proposition:ch04-Chebyshev}
令 $f$ 为 $E$ 上的非负可测函数. 则对于任何 $\lambda>0$,
\begin{align}
m\{x\in E\mid f(x)\geqslant\lambda\}\leqslant\frac{1}{\lambda}\cdot\int_{E}f.\label{align:ch04-Chebyshev}
\end{align}
\end{proposition}

\begin{proof}
定义 $E_{\lambda}=\{x\in E\mid f(x)\geqslant\lambda\}$. 首先假定 $m(E_{\lambda})=\infty$. 令 $n$ 为自然数. 定义 $E_{\lambda,n}=E_{\lambda}\cap[-n,n]$ 与 $\psi_{n}=\lambda\cdot\chi_{E_{\lambda,n}}$. 则 $\psi_{n}$ 是有限支撑的有界可测函数,
\begin{align*}
\text{在 }E\text{ 上对所有 }n,\quad\lambda\cdot m(E_{\lambda,n})=\int_{E}\psi_{n}\quad\text{且 }0\leqslant\psi_{n}\leqslant f.
\end{align*}
我们从测度的连续性推出,
\begin{align*}
\infty=\lambda\cdot m(E_{\lambda})=\lambda\cdot\lim_{n\to\infty}m(E_{\lambda,n})=\lim_{n\to\infty}\int_{E}\psi_{n}\leqslant\int_{E}f.
\end{align*}
由于两边都等于 $\infty$, 因此不等式\eqref{align:ch04-Chebyshev}成立. 现在考虑 $m(E_{\lambda})<\infty$ 的情形. 定义 $h=\lambda\cdot\chi_{E_{\lambda}}$. 则 $h$ 是一个有限支撑的有界可测函数且在 $E$ 上 $0\leqslant h\leqslant f$. 根据 $f$ 在 $E$ 上的积分的定义,
\begin{align*}
\lambda\cdot m(E_{\lambda})=\int_{E}h\leqslant\int_{E}f.
\end{align*}
该不等式两边除以 $\lambda$ 得到 Chebyshev 不等式.
\end{proof}

\begin{proposition}\label{proposition:ch04-积分为零}
令 $f$ 为 $E$ 上的非负可测函数. 则
\begin{align}
\int_{E}f=0\quad\text{当且仅当在 }E\text{ 上 }f=0\text{ a.e.}\label{align:ch04-积分为零}
\end{align}
\end{proposition}

\begin{proof}
首先假设 $\int_{E}f=0$. 则根据 Chebyshev 不等式, 对每个自然数 $n$, $m\{x\in E\mid f(x)\geqslant 1/n\}=0$. 根据 Lebesgue 测度的可数可加性, $m\{x\in E\mid f(x)>0\}=0$. 反之, 假定在 $E$ 上 $f=0$ a.e.. 令 $\varphi$ 为简单函数而 $h$ 为有限支撑的有界可测函数, 使得在 $E$ 上 $0\leqslant\varphi\leqslant h\leqslant f$. 则在 $E$ 上 $\varphi=0$ a.e. 且因此 $\int_{E}\varphi=0$. 由于这对所有这样的 $\varphi$ 成立, 我们推出 $\int_{E}h=0$. 由于这对所有这样的 $h$ 成立, 我们推出 $\int_{E}f=0$.
\end{proof}

\begin{theorem}[积分的线性与单调性]\label{theorem:ch04-非负积分线性}
令 $f$ 和 $g$ 为 $E$ 上的非负可测函数. 则对任何 $\alpha>0$ 和 $\beta>0$,
\begin{align}
\int_{E}(\alpha f+\beta g)=\alpha\int_{E}f+\beta\int_{E}g.\label{align:ch04-非负线性}
\end{align}
此外,
\begin{align}
\text{若在 }E\text{ 上 }f\leqslant g,\quad\text{则}\int_{E}f\leqslant\int_{E}g.\label{align:ch04-非负单调}
\end{align}
\end{theorem}

\begin{proof}
对于 $\alpha>0$, 在 $E$ 上 $0\leqslant h\leqslant f$ 当且仅当在 $E$ 上 $0\leqslant\alpha h\leqslant\alpha f$. 因此, 根据有限支撑的有界函数的积分的线性, $\int_{E}\alpha f=\alpha\int_{E}f$. 因此, 为证明线性我们仅须考虑 $\alpha=\beta=1$ 的情形. 令 $h$ 和 $k$ 为满足在 $E$ 上 $0\leqslant h\leqslant f$ 与 $0\leqslant k\leqslant g$ 的有限支撑的有界可测函数. 我们有在 $E$ 上 $0\leqslant h+k\leqslant f+g$, 而 $h+k$ 也是有限支撑的有界可测函数. 因此, 根据有限支撑的有界可测函数的积分的线性,
\begin{align*}
\int_{E}h+\int_{E}k=\int_{E}(h+k)\leqslant\int_{E}(f+g).
\end{align*}
当 $h$ 和 $k$ 在满足 $h\leqslant f$ 且 $k\leqslant g$ 的有限支撑的有界可测函数中变动时, 左边积分和的最小上界等于 $\int_{E}f+\int_{E}g$. 这些不等式告诉我们, $\int_{E}(f+g)$ 是这些同样和的上界. 因此,
\begin{align*}
\int_{E}f+\int_{E}g\leqslant\int_{E}(f+g).
\end{align*}
剩下来要证明反方向的不等式, 即
\begin{align*}
\int_{E}(f+g)\leqslant\int_{E}f+\int_{E}g.
\end{align*}
当 $\ell$ 在 $E$ 上满足 $0\leqslant\ell\leqslant f+g$ 的有限支撑的有界可测函数中变动时, 根据 $\int_{E}(f+g)$ 作为 $\int_{E}\ell$ 的上确界的定义, 为证明该不等式, 充分且必要的是证明: 对任何这样的函数 $\ell$,
\begin{align}
\int_{E}\ell\leqslant\int_{E}f+\int_{E}g.\label{align:ch04-拆分估计}
\end{align}
对任何这样的函数 $\ell$, 定义 $E$ 上的函数 $h$ 和 $k$ 为
\begin{align*}
\text{在 }E\text{ 上},\quad h=\min\{f,\ell\}\quad\text{与 }k=\ell-h.
\end{align*}
令 $x$ 属于 $E$. 若 $\ell(x)\leqslant f(x)$, 则 $k(x)=0\leqslant g(x)$; 若 $\ell(x)>f(x)$, 则 $k(x)=\ell(x)-f(x)\leqslant g(x)$. 因此, 在 $E$ 上 $k\leqslant g$. $h$ 和 $k$ 都是有限支撑的有界可测函数. 我们有
\begin{align*}
\text{在 }E\text{ 上},\quad 0\leqslant h\leqslant f,\quad 0\leqslant k\leqslant g\quad\text{且 }\ell=h+k.
\end{align*}
因此, 又一次用有限支撑的有界可测函数的积分的线性与 $\int_{E}f$ 和 $\int_{E}g$ 的定义, 我们有
\begin{align*}
\int_{E}\ell=\int_{E}h+\int_{E}k\leqslant\int_{E}f+\int_{E}g.
\end{align*}
因此\eqref{align:ch04-拆分估计}成立, 而线性的证明完成了.

鉴于 $\int_{E}f$ 作为上确界的定义, 为证明单调性不等式\eqref{align:ch04-非负单调}, 充分且必要的是证明: 若 $h$ 是满足在 $E$ 上 $0\leqslant h\leqslant f$ 的有限支撑的有界可测函数, 则
\begin{align}
\int_{E}h\leqslant\int_{E}g.\label{align:ch04-单调证明}
\end{align}
令 $h$ 为这样的一个函数. 则在 $E$ 上, $h\leqslant g$. 因此, 根据 $\int_{E}g$ 作为上确界的定义, $\int_{E}h\leqslant\int_{E}g$. 这完成了单调性的证明.
\end{proof}

\begin{theorem}[积分在区域上的可加性]\label{theorem:ch04-区域可加}
令 $f$ 为 $E$ 上的非负可测函数. 若 $A$ 和 $B$ 是 $E$ 的不交可测子集, 则
\begin{align*}
\int_{A\cup B}f=\int_{A}f+\int_{B}f.
\end{align*}
特别地, 若 $E_{0}$ 是 $E$ 的测度为零的子集, 则
\begin{align}
\int_{E}f=\int_{E\sim E_{0}}f.\label{align:ch04-零测忽略}
\end{align}
\end{theorem}

\begin{proof}
积分在区域上的可加性从线性得到, 如同对有限测度集上的有界函数. 分割公式\eqref{align:ch04-零测忽略}可从区域上的可加性以及由命题\cref{proposition:ch04-积分为零}观察到的非负函数在零测度集上的积分是零得到.
\end{proof}

以下引理使我们能得到几个准则以确保积分号下取极限的合理性.

\begin{lemma}[Fatou 引理]\label{lemma:ch04-Fatou}
令 $\{f_{n}\}$ 为 $E$ 上的非负可测函数序列.
\begin{align}
\text{若在 }E\text{ 上 a.e. 逐点 }\{f_{n}\}\to f,\quad\text{则}\int_{E}f\leqslant\liminf\int_{E}f_{n}.\label{align:ch04-Fatou}
\end{align}
\end{lemma}

\begin{proof}
基于\eqref{align:ch04-零测忽略}, 通过可能地从 $E$ 去掉测度为零的集合, 我们假设在整个 $E$ 逐点收敛. 函数 $f$ 是非负与可测的, 这是由于它是这样的函数序列的逐点极限. 为了证明\eqref{align:ch04-Fatou}中的不等式, 必要且充分的是证明: 若 $h$ 是任何满足在 $E$ 上 $0\leqslant h\leqslant f$ 的有限支撑的有界可测函数, 则
\begin{align}
\int_{E}h\leqslant\liminf\int_{E}f_{n}.\label{align:ch04-Fatou关键}
\end{align}
令 $h$ 为这样的一个函数. 选取 $M\geqslant 0$ 使得在 $E$ 上 $|h|\leqslant M$. 定义 $E_{0}=\{x\in E\mid h(x)\neq 0\}$. 则 $m(E_{0})<\infty$. 令 $n$ 为自然数. 定义 $E$ 上的函数 $h_{n}$ 为
\begin{align*}
\text{在 }E\text{ 上},\quad h_{n}=\min\{h,f_{n}\}.
\end{align*}
观察到函数 $h_{n}$ 是可测的,
\begin{align*}
\text{在 }E_{0}\text{ 上},\quad 0\leqslant h_{n}\leqslant M;\quad\text{而在 }E\sim E_{0}\text{ 上},\quad h_{n}\equiv 0.
\end{align*}
此外, 对 $E$ 中的每个 $x$, 由于 $h(x)\leqslant f(x)$ 且 $\{f_{n}(x)\}\to f(x)$, $\{h_{n}(x)\}\to h(x)$. 从应用于 $h_{n}$ 在有限测度集 $E_{0}$ 上的限制的一致有界序列的有界收敛定理, 以及每个 $h_{n}$ 在 $E\sim E_{0}$ 等于零, 我们推出
\begin{align*}
\lim_{n\to\infty}\int_{E}h_{n}=\lim_{n\to\infty}\int_{E_{0}}h_{n}=\int_{E_{0}}h=\int_{E}h.
\end{align*}
然而, 对每个 $n$, 在 $E$ 上 $h_{n}\leqslant f_{n}$. 因此, 根据 $f_{n}$ 在 $E$ 上的积分的定义, $\int_{E}h_{n}\leqslant\int_{E}f_{n}$. 因此,
\begin{align*}
\int_{E}h=\lim_{n\to\infty}\int_{E}h_{n}\leqslant\liminf\int_{E}f_{n}.
\end{align*}
\end{proof}

Fatou 引理的不等式可以是严格的.

\begin{example}\label{example:ch04-Fatou严格}
令 $E=(0,1]$, 且对自然数 $n$, 定义 $f_{n}=n\cdot\chi_{(0,1/n)}$. 则 $\{f_{n}\}$ 在 $E$ 上逐点收敛于 $f\equiv 0$. 然而,
\begin{align*}
\int_{E}f=0<1=\lim_{n\to\infty}\int_{E}f_{n}.
\end{align*}
Fatou 引理严格不等式的另一个例子: 令 $E=\mathbf{R}$, 且对自然数 $n$, 定义 $g_{n}=\chi_{(n,n+1)}$. 则 $\{g_{n}\}$ 在 $E$ 上逐点收敛于 $g\equiv 0$. 然而,
\begin{align*}
\int_{E}g=0<1=\lim_{n\to\infty}\int_{E}g_{n}.
\end{align*}
然而, 若序列 $\{f_{n}\}$ 是递增的, Fatou 引理中的不等式成为等式.
\end{example}

\begin{theorem}[单调收敛定理]\label{theorem:ch04-MCT}
令 $\{f_{n}\}$ 为 $E$ 上的非负可测函数的递增序列.
\begin{align*}
\text{若在 }E\text{ 上 a.e. 逐点 }\{f_{n}\}\to f,\quad\text{则}\lim_{n\to\infty}\int_{E}f_{n}=\int_{E}f.
\end{align*}
\end{theorem}

\begin{proof}
根据 Fatou 引理,
\begin{align*}
\int_{E}f\leqslant\liminf\int_{E}f_{n}.
\end{align*}
然而, 对每个指标 $n$, 在 $E$ 上 $f_{n}\leqslant f$ a.e., 根据非负可测函数的积分的单调性和\eqref{align:ch04-零测忽略}, $\int_{E}f_{n}\leqslant\int_{E}f$. 因此
\begin{align*}
\limsup\int_{E}f_{n}\leqslant\int_{E}f.
\end{align*}
因此
\begin{align*}
\int_{E}f=\lim_{n\to\infty}\int_{E}f_{n}.
\end{align*}
\end{proof}

\begin{corollary}\label{corollary:ch04-级数积分}
令 $\{u_{n}\}$ 为 $E$ 上的非负可测函数序列. 若在 $E$ 上 a.e. 逐点 $f=\sum_{n=1}^{\infty}u_{n}$, 则 $\int_{E}f=\sum_{n=1}^{\infty}\int_{E}u_{n}$.
\end{corollary}

\begin{proof}
对每个指标 $n$, 应用单调收敛定理于 $f_{n}=\sum_{k=1}^{n}u_{k}$, 接着用非负可测函数的积分的线性.
\end{proof}

\begin{definition}[非负可积函数]\label{definition:ch04-非负可积}
可测集 $E$ 上的非负可测函数 $f$ 称为在 $E$ 上可积, 若
\begin{align*}
\int_{E}f<\infty.
\end{align*}
\end{definition}

\begin{proposition}\label{proposition:ch04-可积有限}
令非负函数 $f$ 在 $E$ 上可积. 则 $f$ 在 $E$ 上 a.e. 有限.
\end{proposition}

\begin{proof}
令 $n$ 为自然数. Chebyshev 不等式和测度的单调性告诉我们
\begin{align*}
m\{x\in E\mid f(x)=\infty\}\leqslant m\{x\in E\mid f(x)\geqslant n\}\leqslant\frac{1}{n}\int_{E}f.
\end{align*}
但 $\int_{E}f$ 是有限的, 因此 $m\{x\in E\mid f(x)=\infty\}=0$.
\end{proof}

\begin{lemma}[Beppo Levi 引理]\label{lemma:ch04-BeppoLevi}
令 $\{f_{n}\}$ 为 $E$ 上的非负可测函数的递增序列. 若积分序列 $\left\{\int_{E}f_{n}\right\}$ 是有界的, 则在 $E$ 上 $\{f_{n}\}$ 逐点收敛于 $E$ 上的 a.e. 有限的可测函数 $f$ 且
\begin{align*}
\lim_{n\to\infty}\int_{E}f_{n}=\int_{E}f<\infty.
\end{align*}
\end{lemma}

\begin{proof}
每个单调的扩充实值函数序列收敛于某个扩充实数. 由于 $\{f_{n}\}$ 是 $E$ 上的递增的扩充实值函数序列, 我们可以在 $E$ 上逐点定义扩充实值非负函数 $f$ 为
\begin{align*}
\text{对所有 }x\in E,\quad f(x)=\lim_{n\to\infty}f_{n}(x).
\end{align*}
根据单调收敛定理, $\left\{\int_{E}f_{n}\right\}\to\int_{E}f$. 因此, 由于实序列 $\left\{\int_{E}f_{n}\right\}$ 是有界的, 它的极限是有限的, 因而 $\int_{E}f<\infty$. 我们从前一个命题推出 $f$ 在 $E$ 上 a.e. 有限.
\end{proof}

\begin{exercise}\label{exercise:ch04-17}
令 $E$ 为测度为零的集合且在 $E$ 上定义 $f\equiv\infty$. 证明 $\int_{E}f=0$.
\end{exercise}

\begin{exercise}\label{exercise:ch04-18}
证明具有有限支撑的有界可测函数的积分是正确定义的.
\end{exercise}

\begin{exercise}\label{exercise:ch04-19}
对于数 $\alpha$, 定义 $f(x)=x^{\alpha}$ ($0<x\leqslant 1$), 而 $f(0)=0$. 计算 $\int_{0}^{1}f$.
\end{exercise}

\begin{exercise}\label{exercise:ch04-20}
令 $\{f_{n}\}$ 为 $E$ 上的逐点收敛于 $f$ 的非负可测函数序列. 令 $M\geqslant 0$, 使得对所有 $n$ $\int_{E}f_{n}\leqslant M$. 证明 $\int_{E}f\leqslant M$. 验证这个性质等价于 Fatou 引理.
\end{exercise}

\begin{exercise}\label{exercise:ch04-21}
令函数 $f$ 在 $E$ 上非负且可积, 而 $\varepsilon>0$. 证明存在 $E$ 上具有有限支撑的简单函数 $\eta$, 在 $E$ 上 $0\leqslant\eta\leqslant f$ 且 $\int_{E}|f-\eta|<\varepsilon$. 若 $E$ 是闭有界区间, 证明存在 $E$ 上的阶梯函数 $h$, 它具有有限支撑且 $\int_{E}|f-h|<\varepsilon$.
\end{exercise}

\begin{exercise}\label{exercise:ch04-22}
令 $\{f_{n}\}$ 为 $\mathbf{R}$ 上的非负可测函数序列, 它在 $\mathbf{R}$ 上逐点收敛于 $f$ 而 $f$ 在 $\mathbf{R}$ 上可积. 证明
\begin{align*}
\text{若}\int_{\mathbf{R}}f=\lim_{n\to\infty}\int_{\mathbf{R}}f_{n},\quad\text{则对任何可测集 }E,\ \int_{E}f=\lim_{n\to\infty}\int_{E}f_{n}.
\end{align*}
\end{exercise}

\begin{exercise}\label{exercise:ch04-23}
令 $\{a_{n}\}$ 为非负实数序列. 定义 $E=[1,\infty)$ 上的函数 $f$ 为: $f(x)=a_{n}$ ($n\leqslant x<n+1$). 证明 $\int_{E}f=\sum_{n=1}^{\infty}a_{n}$.
\end{exercise}

\end{document}